\section{The Hidden Subgroup Problem}

\subsection{Simon's Problem}
Simon's problem is on first glance similar to the problem of Bernstein and Vazirani. Again, we are dealing with an oracle in the form of a function on $N$ bits and a hidden $N$-bitstring $\mathbf{y}$. 
The first difference: the function $f$ is not mapping to a single bit, but can map to many
\begin{align}
    f: \left\{0,1\right\}^N \rightarrow \left\{0,1\right\}^M,\;\; M\leq N
\end{align}
The oracle $f$ promises
\begin{align}
    \forall \mathbf{x} \neq \mathbf{x}': f(\mathbf{x}) = f(\mathbf{x}') \Leftrightarrow \mathbf{x}' = \mathbf{x} \oplus \mathbf{y}
\end{align}
and the task is to find the \textit{hidden} bitstring $\mathbf{y}$.\\

A classical algorithm to the problem would be
\begin{enumerate}
    \item Select a pair $\mathbf{x},\mathbf{x'}$ (at random).
    \item evaluate $f(\mathbf{x})=f(\mathbf{x')}$, if true keep the pair.
    \item repeat until $M$ pairs are found.
    \item The $M$ pairs $\left(\mathbf{x}_k,\mathbf{x}'_k\right)$ provide enough information to solve the linear system $\mathbf{x}_k = \mathbf{x}'_k \oplus \mathbf{y}$ and determine $\mathbf{y}$.
\end{enumerate}
The problem with efficiency is, that, in the worst-case, we need to sample an exponential amount of pairs (leading to exponentially many queries to $f$). The quantum version allows us to find $M$ pairs with a linear amount of queries. An average case analysis improves things classically, but remains exponential, we will not bother with this here.\\

\begin{def*}[Simon's Problem]
We are given a function $f:\left\{0,1\right\}^N \rightarrow \left\{0,1\right\}^M$ with
$$ f(\mathbf{x}) = f(\mathbf{z}) \Leftrightarrow \mathbf{x} \oplus \mathbf{y} = \mathbf{z}$$
for all $\mathbf{x}\neq \mathbf{z}$ and unique hidden $N$-bitstring $\mathbf{y}\in\left\{0,1\right\}^N$ that has to be found.
\end{def*}

We can already recover a special case: If $\mathbf{y}=\mathbf{0}$ then $\mathbf{x} = \mathbf{x}' \oplus \mathbf{y}$ holds true for \textit{all} pairs $(\mathbf{x}, \mathbf{x}')$, therefore $f(\mathbf{x})=f(\mathbf{x}')$ has to hold for all pairs as well making $f$ constant.~\footnote{Note however, that we can not simply run \texttt{Deutsch-Josza} to identify this case, since after a single run that results in a positive answer, we can not be sure, if the probability was 100\%. The Algorithm needs to be adapted a little, and we will require multiple queries to it, in order to recover $\mathbf{y}$.}
On the other hand, if the hidden string is not the zero-string, $\mathbf{y}\neq 0$, then the function needs to have $2^N/2$ distinct values: one for each pair $\mathbf{x}, \mathbf{x}\oplus \mathbf{y}$.~\footnote{proof: for $\mathbf{a} \neq \mathbf{b} \oplus \mathbf{y}$ assume $f(\mathbf{a})=f(\mathbf{b})$ then it follows from the promised structure $\mathbf{a} = \mathbf{b} \oplus \mathbf
y$ which is a contradiction. At the same time, if we have $\mathbf{a} = \mathbf{a}' \oplus \mathbf{y}$ we can not have a second $\mathbf{a} = \mathbf{a}'' \oplus \mathbf{y}$ because  $\mathbf{a}  = \mathbf{a}' \oplus \mathbf{y}= \mathbf{a}'' \oplus \mathbf{y} \Rightarrow \mathbf{a}'=\mathbf{a}''$ }\\
We can summarize the promised properties of the function $f$ as
\begin{align}
&f(\mathbf{x}) = f(\mathbf{x} \oplus \mathbf{y})\\
&\mathbf{y}\neq\mathbf{0} \Rightarrow f(\mathbf{x}) \neq f(\mathbf{z})\; \; \forall \mathbf{z} \notin \left\{\mathbf{x}, \mathbf{x}\oplus \mathbf{y}\right\}\\
&\mathbf{y}=\mathbf{0} \Rightarrow f=\text{const}
\end{align}


The algorithm works with two registers, one carrying the input of $f$ and one the output. We assume that the operation $f$ is available as unitary operation $U_f$ acting on these registers as
\begin{align}
    U_f \left(\ket{\mathbf{x}}\otimes \ket{\mathbf{0}}\right) = \ket{\mathbf{x}}\otimes \ket{ f(\mathbf{x})}.
\end{align}
The full algorithm is
\begin{algo*}[\texttt{Simon}{'s Algorithm}]
Execute the circuit
\begin{center}
$ U_\text{S}$ = \begin{quantikz}
& \gate{\hd^{\otimes N}} & \gate[2]{U_f}& \gate{\hd^{\otimes N}} & \qw \\
& \qw  &[1cm] & \qw & \qw\\
\end{quantikz}
\end{center}
and keep the results on the first register. Repeat $O(N)$ times until enough bitstrings $\mathbf{k}$ are sampled to solve the linear system $\mathbf{y}\odot \mathbf{k}=0$.
\end{algo*}
We begin (at this point, almost traditionally) with a sequence of Hadamard gates on each qubit in the first register followed by the oracle $U_f$ acting on both
\begin{align}
    \ket{\mathbf{0}}\otimes \ket{\mathbf{0}} \xrightarrow[]{\hd^{\otimes N}} \ket{+}^{\otimes N} \otimes \ket{\mathbf{0}} = \frac{1}{\sqrt{2^N}}\left(\sum_{k=0}^{2^N-1}\ket{\mathbf{k}}\right) \otimes \ket{\mathbf{0}} \xrightarrow[]{U_f} \frac{1}{\sqrt{2^N}} \sum_{k=0}^{2^N-1} \left(\vphantom{\sum_{\suum}}\ket{\mathbf{k}}\otimes \ket{{f(\mathbf{k})}}\right).\nonumber
\end{align}
After the state is prepared, we read out the second qubit register -- that is, we will get one of the possible values of $f(\mathbf{k})$, let's call them $a = f(\mathbf{i})$, with $\mathbf{i}\in \left\{\mathbf{k}\right\}_{k=0}^{2^N-1}$ and probability of $1/2^N$ for each of them. As a result the remaining qubit register collapses to a linear combination of $\mathbf{i}$ and the partner $\mathbf{i}\oplus \mathbf{y}$ that lead to $f(\mathbf{i})=a$
\begin{align}
    \frac{1}{\sqrt{2}} \left( \ket{\mathbf{i}} + \ket{\mathbf{i\oplus y}}\right),
\end{align}
At this point we do not profit yet, from reading out the first register, we first need to transform again with a Hadamard operation on all qubits on the remaining register
\begin{align}
        \frac{1}{\sqrt{2}} \left( \ket{\mathbf{i}} + \ket{\mathbf{i \oplus y}}\right) \xrightarrow[]{\hd^{\otimes N}}& \frac{1}{\sqrt{2}}\left( \left(\frac{1}{\sqrt{2^N}}\sum_{k=0}^{2^N-1} (-1)^{\mathbf{k}\odot \mathbf{i}}\ket{\mathbf{k}}\right) +\left(\frac{1}{\sqrt{2^N}}\sum_{k=0}^{2^N-1} (-1)^{\mathbf{k}\odot( \mathbf{i}\oplus\mathbf{y})}\ket{\mathbf{k}}\right) \right)\nonumber\\
        &= \frac{1}{\sqrt{2^{N+1}}}\sum_{k=0}^{2^N-1} (-1)^{\mathbf{k\odot i}} \left(1+(-1)^{\mathbf{k}\oplus \mathbf{y}}\right) \ket{\mathbf{k}} \nonumber\\
        & = \frac{1}{\sqrt{|A|}} \sum_{k\in A^\perp} (-1)^{\mathbf{k\odot i}} \ket{\mathbf{k}}, \quad A^\perp = \left\{\mathbf{k} \;|\; \mathbf{k} \odot \mathbf{y} = 0\right\}.
\end{align}
This is where the magic happened. If we are now reading out the register, we will get a $\mathbf{k}$ such that $\mathbf{k\odot \mathbf{y}}=0$ (with uniform probability over all such candidates). This is completely independent from what value $a=f(\mathbf{i})$ we have read out before. To determine the $M$-bitstring $\mathbf{y}$ via the linear system $\mathbf{k\odot \mathbf{y}}=0$ it suffices to get $M$ candidates $\mathbf{k}$. So we need to repeat the algorithm $O(M)$ times to solve the problem. While doing this, we can also detect the special case $f=const$ with high probability (we would read out the same value in the first register on every trial and the second register would always give back $\mathbf{0}$).
To digest this information, it is beneficial to look at a concrete example.

\paragraph{Example: $N=2,M=1,\; f(\mathbf{x})=x_1\oplus x_2$}: In this example our function acts on two qubits with the concrete assignments $00\rightarrow 0, 01\rightarrow 1, 10 \rightarrow 1, 11 \rightarrow0$. So we have $f(\mathbf{0})=f(\mathbf{3})$ and $f(\mathbf{1})=f(\mathbf{2})$ and the hidden string is $\mathbf{y}=\mathbf{3}=11$. From all possible 4 combination of two bits, two (00 and 11) fulfill the equation $\mathbf{y} \odot \mathbf{k} = y_1k_1 \oplus y_2k_2 = k_1 \oplus k_2=0$. The quantum circuit prepares the states
\begin{align}
    \ket{0}^{\otimes}\otimes \ket{0} \xrightarrow[]{\hd^{\otimes N}}& \frac{1}{2}\left(\ket{00} + \ket{01} + \ket{10} + \ket{11}\right)\otimes \ket{0}\\
    \xrightarrow[]{U_f}& \frac{1}{2} \left(\ket{00}\otimes\ket{0} + \ket{01}\otimes\ket{1} + \ket{10}\otimes \ket{1} + \ket{11}\otimes \ket{0}\right)\nonumber\\
    \xrightarrow[]{\hd^{\otimes N}}& \frac{1}{2} \left(\ket{++}\otimes\ket{0} + \ket{+-}\otimes\ket{1} + \ket{-+}\otimes \ket{1} + \ket{--}\otimes \ket{0}\right)\nonumber\\
    &= \frac{1}{\sqrt{2}}\left(\left(\ket{00} + \ket{11}\right)\otimes\ket{0} + \left(\ket{00} + \ket{11}\right)\otimes\ket{1}\right)\nonumber\\
    & =\frac{1}{\sqrt{2}}\left(\ket{00} + \ket{11}\right)\otimes\ket{+}.
\end{align}
So by reading out the first two qubits we either get $00$ or $11$, the two candidates we need to determine $\mathbf{y}$ via the linear system $\mathbf{y}\odot \mathbf{k}=0$. Note again that the end result is independent of the third qubit.
At this point we also realize, that reading out the second register (that holds the value of $f$ after applying the oracle) is actually not necessary. 
\subsection{Groups and their Properties}

A group is an algebraic structure that obeys the axioms
\begin{def*}[(Abelian) Group $(G,\odot)$]\phantom{...}\\
A set of elements $G = \left\{e_k\right\}_{k=1}^N$ with group action $\odot: G\times G
\rightarrow G $ and axioms
\begin{align*}
    &a,b,c\in G: (a\odot b)\odot c = a\odot (b\odot c)& \text{associativity}\\
    &\exists e\in G, \forall g\in G:  e\odot g = g\odot e = g &\text{neutral element}\\
    &\forall g\in G, \exists g^{-1}\in G: g\odot g^{-1}= g^{-1}\odot g = e & \text{inverse element}
\end{align*}
If additionally the group action commutes
\begin{align}
    \forall a,b\in G: a \odot b = b\odot a
\end{align}
the group is called \textit{abelian}
\end{def*}
Note that we used $\odot$ as a group action and not as a binary inner product to follow standard notation. If we are dealing with a concrete realization or a \textit{representation} of a group, we replace $\odot$ with the corresponding symbol that is most suitable (e.g. $\odot \rightarrow +$ if we are using addition of numbers.\\
The usual examples are: 
\begin{itemize}
    \item $(\mathbb{Z},+)$ (integers under addition)
    \item $(\mathbb{C}^{N\times M}, +)$ (matrices under addition)
    \item $(\text{Herm}(N),+)$ (Hermitian $N\times N$ matrices under addition)
    \item $(U(N), \cdot )$ (Unitary $N\times N$ matrices under matrix-multiplication)
\end{itemize}
which are, except for the last one, all abelian.
To test your understanding of the notion of a group: show that $(\mathbb{Z},\cdot)$, $(\mathbb{R}, \cdot)$ and $(\mathbb{N}, +)$ are not groups and explain why general matrices under matrix multiplication do not form a group.\\

Prominent Abelian groups are the cyclic groups that can be represented by integers with modular addition or rotations in vector spaces.
\begin{def*}[Cyclic Group]
The cyclic group $C_n$ is an abelian group of $n$ elements $C_n = \left\{g_0,g_1,g_2,g_3,\dots g_{n-1}\right\}$ with
\begin{align}
    &g_0\odot g_k = g_k & \text{neutral element}\\
    &g_1 \odot g_k = g_{(k\oplus_n1) } & \text{cyclic property}
\end{align}
where $\oplus_n$ denotes addition modulo $n$.
The element $g_1$ is called the \textit{generator} of the cyclic group and a typical notation is $G=\left\{e,g,g^2,g^3,\dots, g^{n-1}\right\}$.
\end{def*}
This is indeed a group: the inverse element of $g^k$ is $g^{n-k}$, since the neutral element is $e=g^0=g^n$. The typical example of a representation of $C_n$ is $\mathbb{Z}/n\mathbb{Z}$ integers with addition modulo $n$. 
A single bit is a realization of the group $C_2$ with $(\left\{0,1\right\},\oplus)$, two bits form the group $\left(\left\{0,1,2,3\right\}, \oplus_3\right)$. 
If we use bitwise binary addition, denoted as a simple $\oplus$, we have further examples for groups: $(\left\{\mathbf{0},\mathbf{1},\mathbf{2},\mathbf{3}\right\}, \oplus)$~\footnote{e.g. $\mathbf{1}\oplus\mathbf{3} = (0,1)\oplus(1,1) = (0\oplus 1, 1\oplus 1) = (1,0) = \mathbf{2}$} or more general $(\left\{\mathbf{k}\right\}_{k=0}^{2^N-1}, \oplus)$ we have a group with neutral element $\mathbf{0}$ and self-inverse elements. On first glance this doesn't appear to be a cyclic group, and it isn't, but it is isomorphic to the direct product of $N$ cyclic $C_2$ groups: $C^N_2:=C_2\times C_2\times \dots \times C_2$.\\
To continue with the quantum algorithms we unfortunately still need some vocabulary.

\begin{def*}[Subgroup]\phantom{...}\\
A subgroup $A\leq G$ of a group $(G,\odot)$ is a subset $A\subseteq G$ of the group elements such that $(A,\odot)$ is a group. Trivial subgroups are $A=\left\{e\right\}$ and $A=G$. The 
\end{def*}
Examples: The groups $(\mathbb{Z}, +)$ has no non-trivial subgroups. The group $(\text{Herm}(N), +)$ is a subgroup of $(\mathbb{C}^{N\times M}, +)$. The group $SU(N): \left\{e\in U(N)\;|\; \det(e)=1\right\}$ is a subgroup of $U(N)$ and each set of bitstrings with $M\leq N$ set to $0$ is a subgroup of $C_2^N$.\
\begin{def*}[Generating Set]\phantom{...}\\
A generating set $S\subseteq G$ of a group $(G,\odot)$ has the property
\begin{align*}
    \forall g\in G:  g=s_1\odot s_2\odot \dots \odot s_M
\end{align*}
with finite $M$ and $s_k\in S$ or $s_k^{-1}\in S$.
Meaning every element of the group can be generated by a finite combination of generating elements under the group action. If $S$ generates $G$ we write $G = \langle S\rangle $.
\end{def*}
Example: The notion of a generating set gives us an alternative definition of cyclic groups: the groups that can be generated by a single element $C_n = \langle \left\{g\right\}\rangle $ but also $C_n = \langle g^2\rangle$ and so on. The group of integers under addition  $(\mathbb{Z}, +)=\langle \left\{1\right\}\rangle $ is generated by a single element.~\footnote{the neutral element is for example generated via $0=1-1$, since the inverse of the generating elements can be used as well. In the same way, we can generate the negative numbers.} A generating set for the group $(\mathbb{R}, +)$ is the set $S=(0,1)\subset \mathbb{R}$ since for any  $x\in \mathbb{R}$ we can chose a $n\in \mathbb{N}$ such that $x/n<1$ meaning that $s=x/n\in S$ if we add this number $n$ times, we get $x$ back: $x=nx/n=ns=s+s+s+\dots +s$. This construction works for every $x\in \mathbb{R}$ and there are always finitely many additions: so we have a generating set.
\begin{def*}[Annihilator]\phantom{...}\\
The annihilator of a subgroup $A\leq G$ of an abelian group $G$ is defined as
\begin{align}
    A^\perp := \left\{g\in G \mid \forall a\in A :\chi_g(a)=1\right\}.
\end{align}    
\end{def*}
The annihilator~\footnote{naming is horrible in this context ... sorry} collects all elements whose characters are 1 for the operations in the respective subgroup. The annihilator itself is also a subgroup $A\leq G \Rightarrow A^\perp \leq G$.~\footnote{this is \textit{not} the case for non-abelian groups which also is one of the reasons why we are restricting ourselves here to abelian groups.}

\subsection{Fourier Transforms on Groups}

\begin{def*}[Representation of a Group] \phantom{...}\\
A map $f: G \rightarrow GL(V)$ that maps elements of a group $(G,\odot)$ to linear operators on the vector space $V$ such that the group action is preserved under the linear operators
\begin{align*}
    &c=a\odot b \Rightarrow f(c)=f(a)f(b) 
\end{align*}
is called a representation of the group $G$ on $V$. 
\end{def*}
In other words, the representation is a homomorphism~\footnote{homo (same) morph (form): this essentially means structure preserving} of the group onto linear operators of a vector space. Note that the preservation of the group action also preserves neutral and inverse elements
\begin{align*}
f(a^{-1}) &= f(a)^{-1}\\
    f(e) &= 1.\\
\end{align*}
Some examples:
The cyclic group $C_4$ is fully characterized over the abstract elements $(\left\{e,g,g^2,g^3\right\}, \odot)$ where $g^3 = g\odot g \odot g$ and $g^3\odot g=e$. If we write $C_4=(\left\{0,1,2,3\right\},\oplus_4)$ its not a representation in the sense of the definition, but rather a realization of the group with numbers (we did not map it to linear operators). A one-dimensional representation $f:C_4 \rightarrow GL(\mathbb{C})$ can be realized as mapping to the one-dimensional matrices (which are just numbers) acting on the one-dimensional vector space of complex numbers
\begin{align}
    \left\{e,g,g^2,g^3\right\} \mapsto \left\{1,i,-1,-i\right\}
\end{align}
but also 
\begin{align}
    \left\{e,g,g^2,g^3\right\} \mapsto \left\{1,-1,1,-1\right\}
\end{align}
would be a valid representation. We can also map all elements to $f(e)=f(g)=f(g^2)=g(g^3)=1$ to the unit operation which gives us a trivial representation of the group. Another example that uses two dimensional complex matrices
\begin{align}
f(e) = \begin{pmatrix}
    1 & 0 \\ 0 & 1 \end{pmatrix}, \; f(g) = \begin{pmatrix}
    0 & 1 \\ 1 & 0 \end{pmatrix},\; f(g^2)=\begin{pmatrix}
    1 & 0 \\ 0 & 1 \end{pmatrix}, \; f(g^3)=\begin{pmatrix}
    0 & 1 \\ 1 & 0 \end{pmatrix}
\end{align}
that is also a valid representation, but seems a bit redundant. And indeed we can \textit{reduce} this representation to a direct sum of two one-dimensional representations. To see this, we apply the transformation $M_k = Uf(g^k)U^\dagger$ to all four matrices and we use the matrix~\footnote{You might have recognized that this is the Hadamard matrix, but this is a coincidence (or bias of the lecturer when thinking about examples)}
\begin{align}
    U = \frac{1}{\sqrt{2}}\begin{pmatrix}
        1 & 1\\1 & -1
    \end{pmatrix}.
\end{align}
The transformed matrices $A_g$ are a representation of the group as well
\begin{align}
    M_0 = \begin{pmatrix} 1 & 0 \\ 0 & 1 \end{pmatrix},\;
    M_1 = \begin{pmatrix} 1 & 0 \\ 0 & -1 \end{pmatrix},\;
    M_2 = \begin{pmatrix} 1 & 0 \\ 0 & 1 \end{pmatrix},\;
    M_3 = \begin{pmatrix} 1 & 0 \\ 0 & -1 \end{pmatrix},\;
\end{align}
since this is diagonal we can treat the first and second components individually, the first components are the trivial representation ($f_1(g^k)=1$) while the second components are the $f_2(g^k)=(-1)^k$ representation we saw above: mathematically this is a direct sum of two one dimensional representations, often written as $f_1 \oplus f_2$. 
For representations of Abelian groups a diagonalizing matrix like the one in the example above can always be found. The individual one-dimensional representations are called \textit{irreducible} representations or short ``irreps''. The individual numbers that make up the one-dimensional representations are called the characters,~\footnote{in general, the character of a representation of a group element is the trace of the corresponding matrix} The group $C_4$ has four irreducible representations
\begin{align}
    A_0=\left\{1,1,1,1\right\},\;
    A_1=\left\{1,-1,1,-1\right\},\; 
    A_2=\left\{1,i,-1,-i\right\},\;
    A_3=\left\{1,-i,-1,i\right\},\; 
\end{align}
three of them we have already encountered. We will denote the characters of them by $\chi_k(g)$ where $k$ enumerates the irreducible representations, e.g. $\chi_1(g)=1$, $\chi_k(e)=1$, $\chi_4(g^3)=-i$. 
\begin{theorem*}[irrep properties]
Irreducible representations of finite Abelian groups
\begin{itemize}
    \item one-dimensional representations (they are $|G|$ complex numbers $e^{i\varphi}$)
    \item there are $|G|$ distinct irresducible representations 
    \item if represented as $|G|$ dimensional vectors they span an orthogonal basis 
\end{itemize}
\end{theorem*}
We won't proof those properties here, but it is a good exercise to at least check if they hold for the examples we encountered.
The last property essentially states that characters of irreducible representation are orthogonal, meaning that we have something like an inner product
\begin{align}
    \sum_{g\in G} \chi_k^*(g) \chi_l(g) = |G|\delta_{kl}.
\end{align}
If we view the characters as normalized vectors, $$\ket{\chi_k} := 1/|G| \begin{pmatrix}\chi_k(e) & \chi_k(g) & \chi_k(g^2) & \dots &\chi_k({g^{|G|-1}}))\end{pmatrix}^T,$$ for example in the case of $C_4$
\begin{align}
    \ket{\chi_1} = \frac{1}{2} \begin{pmatrix}1\\1\\1\\1\end{pmatrix},\;
    \ket{\chi_2} = \frac{1}{2} \begin{pmatrix}1\\-1\\1\\-1\end{pmatrix},\;
    \ket{\chi_3} = \frac{1}{2} \begin{pmatrix}1\\i\\-1\\-i\end{pmatrix},\; 
    \ket{\chi_4} = \frac{1}{2} \begin{pmatrix}1\\-i\\-1\\i\end{pmatrix}, 
\end{align}
we note that they span a basis for $\mathbb{C}^{|G|}$ where $|G|$ denotes the number of elements in the group (e.g. $|C_4|=4$). If we want to transform from the standard basis to the basis spanned by the irreps of a group we can perform the following operation 
\begin{align*}
    U_G\ket{k} = \ket{\chi_k}
\end{align*}
that can be constructed as
\begin{align*}
    U_G = \sum_i \ket{\chi_i}\bra{i}
\end{align*}
with matrix elements
\begin{align}
    \bra{l}U_G\ket{k} = \sum_i \braket{l}{\chi_i}\braket{i}{k} = \chi_k(g_l),
\end{align}
which is also called the \textit{Fourier transform} under the group $G$.
For general cyclic groups $C_N$ the Fourier coefficients (the characters of the group) take a familiar form
\begin{align}
    G=C_N \;:\;\chi_k(g_l) = e^{2\pi i \frac{kl}{N}}
\end{align}
of the ``usual'' discrete Fourier transform. We can recover two other familiar cases
\begin{align}
    C_2 \;:\; \chi_k(g_l) = (-1)^{kl} 
\end{align}
which are the elements of the Hadamard matrix, and\footnote{here $\odot$ denotes the bitwise binary inner product: $\mathbf{k}\odot \mathbf{l} = \left( \sum_{i=1}^{N} k_i l_i \right)\mod 2$}
\begin{align}
    C_{2}^N \;:\;\chi_{k}(g_l) = (-1)^{\mathbf{k}\odot \mathbf{l}}
\end{align}
which are the matrix elements of $\hd\otimes\hd\otimes\dots\otimes\hd$.



\subsection{The Hidden Subgroup Problem}
The quantum algorithm for \texttt{Simon} can be generalized to the hidden subgroup problem (HSP) that aims to find a subgroup $A\leq G$ within a given abelian group $G$.
\begin{def*}[Hidden Subgroup (HSP)]\phantom{...}\\
Consider an abelian group $G$ and a integer valued function $f:G\rightarrow I\subset \mathbb{N}$ that \textit{hides} a subgroup $A\leq G$ through
\begin{align}
    f(x)=f(y) \Leftrightarrow Ax=Ay 
\end{align}
for $x,y\in G$ and with the \textit{cosets}
\begin{align}
    Ax = \left\{a\odot x \;|\; a\in A\right\}.
\end{align}
\end{def*}Now, what does this have to do with Simon's problem? Lets translate:
In Simon's problem, the group function is defined on $N$-bits. The corresponding bitstrings $\mathbf{x}\in\left\{0,1\right\}^N$ form a group -- there could be different group actions, but we can chose bitwise addition $\mathbf{a}\oplus\mathbf{b}:=(a_1,a_2\dots a_N)\oplus(b_1,b_2,\dots b_N)=(a_k\oplus b_k)_{k=1}^N = \mathbf{c}$. The abstract group for this is the direct product of cyclic groups we have encountered in the last section
\begin{align}
    C^N_{2} = C_2 \times C_2 \times C_2 \times \dots \times C_2.
\end{align}
The subgroup to Simon's problem hides the sought bitstring $\mathbf{y}$~\footnote{we need the neutral element $\mathbf{0}$ so that $A$ is a group. One could also rephrase Simon's problem as decision problem that asks if the hidden string is $\mathbf{0}$ or something else.} and is therefore 
\begin{align}
    A=\left\{\mathbf{0},\mathbf{y}\right\}.
\end{align}
Note that $1\leq |A| \leq 2$ since it can happen that $\mathbf{y}=\mathbf{0}$ (this is the special case, we discussed briefly above). 
The annihilator of this group are all elements of $C^N_2$ whose characters of the operations in $A$ are 1, the characters of $C^N_2$ are $\chi_\mathbf{k}(\mathbf{l})= (-1)^{\mathbf{k}\odot \mathbf{l}}$, so the annihilator of the hidden subgroup $A$ is 
\begin{align}
    A^\perp = \left\{\mathbf{k} \;|\; \mathbf{k}\odot \mathbf{y} = 0\right\}
\end{align}
with the bitwise inner product $\odot$.~\footnote{now, the naming of this specific set from before probably makes sense in hindsight, and the name ``annihilator'' as well: what gets annihilated is the generator of the character. }
% First, see that this is indeed a subgroup of $C^N_2$: the neutral element is $\mathbf{0}$ and since $\mathbf{0}\odot \mathbf{y}=0$ for any choice of $\mathbf{y}$ we have $\mathbf{0}\in A$. At the same time, since all members of $C_2^N$ are self inverse, all members of $A$ also are their own inverses. At last we need to verify closedness under the group action (which is bitwise binary addition), suppose we have two $\mathbf{i},\mathbf{j}\in A$, then we can build $\mathbf{k}=\mathbf{i}\oplus \mathbf{j}$ but then $\mathbf{k}\odot\mathbf{y} = \mathbf{i}\odot \mathbf{y} \oplus \mathbf{j}\odot \mathbf{y} = 0\oplus 0$ and therefore $\mathbf{k}\in A$. We have $A\leq C^N_2$, but is it the right subgroup for the hidden subgroup problem?
% Yes, because in Simon's problem we have
% \begin{align}
%     f(\mathbf{x})=f(\mathbf{x\oplus \mathbf{y}})
% \end{align}
% for all $\mathbf{x}\in C^N_2$. So if we have two $\mathbf{i},\mathbf{j}$ for which $f(\mathbf{i}) = f(\mathbf{j})$ then $\mathbf{i} = \mathbf{j} \oplus \mathbf{y}$. The relevant cosets are
% \begin{align}
%     \mathbf{i}A &= \left\{\mathbf{k}\oplus \mathbf{i} \;|\; \mathbf{k}\in A\right\}\nonumber \\
%     &= \left\{\mathbf{k}\oplus (\mathbf{j}\oplus \mathbf{y}) \;|\; \mathbf{k}\in A\right\} \nonumber\\
%     &= \left\{\mathbf{j}\oplus (\mathbf{k}\oplus \mathbf{y}) \;|\; \mathbf{k}\in A\right\} \nonumber\\
%     &= \left\{\mathbf{j}\oplus \mathbf{k}' \;|\; \mathbf{k}'\in A\right\} = \mathbf{j}A
% \end{align}
% where we have used $\mathbf{y}\in A$ and through the group's closedness under its action $\mathbf{k}'=(\mathbf{k}\oplus \mathbf{y})\in A$ for all $\mathbf{k}\in A$.
\subsection{Quantum Algorithm for \texttt{HSP}}
Now that we have a nice guiding example through \texttt{Simon} at hand, let's see how the general quantum algorithm for the \texttt{HSP} looks like. 
We make use of an oracle again, that implements us the function $f$ on two qubit registers
\begin{align}
    U_f\left(\ket{g}\otimes\ket
    0\right)=  \ket{g}\otimes \ket{f(g)}
\end{align}
where the states $\ket{g}$ are encoded into a register with enough qubits to represent all $|G|$ elements of the group of interest. For the example $C^N_2$ these are $N$-qubits and for the $C_4$ example two qubits will be enough. For the second register a single qubit suffices as $f:G\rightarrow \left\{0,1\right\}$ maps to a single bit.
\begin{algo*}[\texttt{HSP}]
Execute the circuit
\begin{center}
$ U_\text{S}$ = \begin{quantikz}
& \gate{\hd^{\otimes N}} & \gate[2]{U_f}& \gate{\mathcal{F}^\dagger} & \qw \\
& \qw  &[1cm] & \qw & \qw\\
\end{quantikz}
\end{center}
\end{algo*}
The first part of the circuit creates the state
\begin{align}
    \ket{0}^{\otimes N} \xrightarrow[]{U_f \hd^{\otimes N}}& \frac{1}{\sqrt{|G|}}\sum_{g\in G}\ket{g}\otimes\ket{f(g)} \nonumber
\end{align}
Read out the second register and get one value $f(h)$ back, the first register is now in superposition of all states $\ket{g}$ such that $f(g)=f(h)$
\begin{align}
 \frac{1}{\sqrt{|B|}}\sum_{b\in B} \ket{b},\;\;B=\left\{g\in G\;|\; f(g)=f(h)\right\}.
\end{align}
From the structural promise of the problem, we know however that these musst be the cosets $bA$ that should be identical for all $b\in B$, so we can just pick the candidate $h$ that we measured and we know $bA=hA$ for all $b\in B$. 
With this insight, we can write the remaining state as
\begin{align}
 \frac{1}{\sqrt{|B|}}\sum_{b\in B} \ket{b} = \frac{1}{\sqrt{|hA|}}\sum_{x\in hA} \ket{x} = \frac{1}{\sqrt{|A|}} \sum_{a\in A} \ket{a\odot h}.
\end{align}
The last part follows from the definition of the coset $hA = \left\{a\odot h \;|\; a\in A\right\}$ -- note that $|A| = |bA|$ holds in general.
We are now at the same stage as in Eq.~\eqref{eq:simon-1}: if we read out the remaining register, we get some information of the unknown element $a$, but it will not really help us yet with regards to efficiency. Similar to \texttt{Simon} we are continuing with a Fourier transform on the group 
\begin{align}
    \ket{\psi}=\frac{1}{\sqrt{|A|}} \sum_{a\in A} \ket{a\cdot h} \xrightarrow[]{\mathcal{F}} \ket{\hat{\psi}} &=\frac{1}{\sqrt{|G|}}\sum_{k\in G} \sum_{a\in A} \chi_k(a\cdot h) \ket{{k}} \nonumber\\
    &= \frac{1}{\sqrt{|G|}} \sum_{k\in G} \sum_{a\in A} \chi_k(a)\chi(h) \ket{k}\\
\end{align}
Now we can use properties of the annihilator
\begin{align}
    \sum_{a\in A}\chi_k(a) = \begin{cases}
        |A|, & k\in A^\perp \\
        0 , & \text{else}
    \end{cases}
\end{align}
and end up with the state
\begin{align}
    \ket{\hat{\psi}}=\frac{|A|}{\sqrt{|G|}} \sum_{k\in A^\perp} \chi_k(h) \ket{k}\\
\end{align}
%which maps the classical basis states (the bitstrings $\ket{\mathbf{k}}$) to the characters of the group. Before, the characters where products of $\ket{\pm}$ states, now they are more abstract states $\ket{\chi_k}$, and depend on the group $G$ of the chosen instance of the problem. The basis states $\ket{a\cdot h}$ are therefore mapped to the corresponding character states $\ket{\chi_{h\cdot A}}$ which are linear combinations of the classical basis states with the group characters as probability amplitudes. At this point it is important to remember, that for characters (which are essentially a one-dimensional representation of the group) it holds $\chi_{h\cdot a} = \chi_{h}\chi_a$. 
In summary the Fourier transform acts as
\begin{align}
     \frac{1}{\sqrt{|A|}} \sum_{a\in A} \ket{a\cdot h} \xrightarrow[]{\mathcal{F}} &\frac{1}{\sqrt{|A|}} \sum_{a\in A} \ket{\chi_{a\cdot h}} = \frac{1}{\sqrt{|A||G|}} \sum_{a\in A} \sum_{g\in G}\chi_{a\cdot h}(g)\ket{g}\nonumber\\
     &= \frac{1}{\sqrt{|A||G|}} \sum_{g\in G} \chi_h(g)\left(\sum_{a\in A}\chi_{a}(g)\right)\ket{g}.
\end{align}
At this point, we need a property of the characters and the definition for generating sets. If the element $g$ is in a generating set of the subgroup $A$, then we have
\begin{align}
    \sum_{a\in A} \chi_a(g) = |A|
\end{align}
and otherwise the sum equates to zero. 
If we use this property for the formulat above we arrive at
\begin{align}
    \frac{1}{\sqrt{|A|}} \sum_{a\in A} \ket{a\cdot h} \xrightarrow[]{\mathcal{F}} & \frac{\sqrt{|A|}}{\sqrt{|G|}} \sum_{g\in \bar{A}} \chi_h(g)\ket{g}
\end{align}
which is a equally weighted (w.r.t to probabilities, because the characters can only be $\pm 1$) superposition of all elements of the generating set of the subgroup. If we read-out the remaining register, we are therefore always getting one element of the generating set. To determine the subgroup $A$ we can therefore expect to get away with $O(|A|)$ repetitions.

\subsection{Examples for the Abelian HSP}